\section{Related Work}
\label{sec:related_work}

Our work resides at the intersection of parameter-space model merging, test-time adaptation, and dynamic neural network architectures. We position QWS-Merge within this landscape below.

\subsection{Parameter-Space Model Merging}
Model merging seeks to combine multiple neural networks with different capabilities into a single unified network without expensive retraining. Early techniques such as weight averaging \cite{frankle2020linear} and permutation alignment \cite{singh2020model, gitrebasin} laid the foundation for merging models starting from the same or different initializations. More recently, Task Arithmetic \cite{ilharco2022editing} demonstrated that task-specific fine-tuned weights can be treated as vectors (task vectors) and added or subtracted to modify model behavior. However, simple linear addition of task vectors often introduces destructive parameter interference. To address this, TIES-Merging \cite{yadav2023ties} resolves parameter sign conflicts and prunes redundant values, while DARE \cite{yu2023language} uses extreme random pruning with rescaling. 

Subsequent work explored regularizing the merging space. For example, RegCalMerge \cite{regcalmerge} utilizes calibration and spatial regularization to mitigate transductive overfitting, while PolyMerge and SplineMerge \cite{polymerge} project coefficients into continuous low-degree polynomial or spline subspaces to filter out high-frequency optimization noise. SAIM \cite{saim} analyzes isotropic sharpness-aware landscapes to ensure flatter minima during merging. Despite these advances, all these methods ultimately produce a \emph{static} merged model, which inevitably suffers from catastrophic representational collapse when task conflicts are high and model capacity is low.

\subsection{Test-Time Adaptation and Few-Shot Parameter Search}
Rather than using fixed, heuristic merging coefficients, test-time adaptation (TTA) search methods have emerged to optimize coefficients on-the-fly. AdaMerging \cite{liang2023adamerging} optimizes layer-wise or task-wise merging coefficients at test-time by minimizing the entropy of the predictions on the test stream without labels. While elegant, unsupervised online entropy minimization is prone to severe transductive overfitting and representation drift, particularly under extreme domain shifts. 

To bypass this ``no-data'' assumption of TTA, OFS-Tune \cite{ofstune} leverages tiny offline labeled validation sets (e.g., 5-10 samples per task) to optimize static, layer-wise merging coefficients under polynomial constraints. ZipMerge \cite{zipmerge} co-optimizes pruning boundaries and merging coefficients on edge devices but identifies the Overfitting-Optimizer Paradox under extreme high-conflict domains—where over-optimizing coefficients on a small validation set degrades out-of-distribution generalization. In contrast, QWS-Merge utilizes a highly constrained, low-dimensional phase-space (only 336 parameters) to model dynamic input-dependent wave functions. This extremely small parameter footprint guarantees stable optimization and robust generalization on small validation sets, completely bypassing the Overfitting-Optimizer Paradox.

\subsection{Dynamic Neural Networks and Wave-Inspired ML}
Dynamic neural networks adapt their architectures or parameters conditioned on the input on-the-fly, providing high representational capacity with low average computational cost. Mixture-of-Experts (MoE) \cite{shazeer2017outrageously, fedus2022switch} routes input tokens to different expert sub-networks. However, traditional MoE requires scaling the total model size, whereas our work aims to merge multiple experts into a \emph{single, fixed-size} backbone.

FoldMerge \cite{foldmerge} proposes a non-linear weight-space diffeomorphism using normalizing flows to warp parameters into a latent coordinate system before blending, but requires complex neural flow matching. Drawing inspiration from physical and quantum mechanics, wave-inspired and quantum-inspired deep learning has been explored in quantum neural networks (QNN) \cite{beer2020training} and Fourier-based network architectures \cite{guibas2021adaptive}. Our method, QWS-Merge, stands apart by porting the physical principle of wavefunction superposition and collapse to parameter space, using wave-like phase interference to dynamically assemble multi-task models at runtime without the computational overhead of non-linear flows.
